\documentclass[11pt]{article}
\usepackage[margin=1in]{geometry}
\usepackage{amsmath}
\usepackage{amssymb}
\usepackage{tcolorbox}
\usepackage{enumitem}

\title{ECO 310: Complete Guide to WARP Questions on Midterms}
\author{All Question Types and How to Solve Them}
\date{March 1, 2026}

\begin{document}

\maketitle

\tableofcontents

\newpage

\section{Overview}

WARP (Weak Axiom of Revealed Preference) appears on almost every ECO 310 midterm, but in different formats. This guide covers ALL the types you might see.

\subsection{What is WARP?}

\textbf{WARP says:} If you choose bundle A when bundle B was also affordable, you cannot later choose bundle B when bundle A is also affordable.

\textbf{Why?} The first choice revealed you prefer A to B. The second choice reveals you prefer B to A. This is a contradiction!

\subsection{The Two Key Conditions for WARP Violation}

For a new bundle $(x_1, x_2)$ to violate WARP when the original bundle was $(x_1^*, x_2^*)$:

\begin{enumerate}
\item \textbf{Condition 1 (Past):} The new bundle was affordable at old prices
\[
p_1^{old} x_1 + p_2^{old} x_2 \leq m^{old}
\]

\item \textbf{Condition 2 (Present):} The old bundle is affordable at new prices (proven by new bundle costing enough)
\[
p_1^{new} x_1 + p_2^{new} x_2 \geq p_1^{new} x_1^* + p_2^{new} x_2^*
\]
\end{enumerate}

\textbf{If BOTH hold $\Rightarrow$ WARP VIOLATION!}

\newpage

\section{Type 1: Given Costs, Determine if WARP is Violated}

\subsection{Format}

You're told:
\begin{itemize}
\item What someone bought on Day 1 and how much it cost
\item What they bought on Day 2 and how much it cost
\item How much one bundle would have cost at the other day's prices
\end{itemize}

\textbf{Question:} Does this violate WARP?

\subsection{Example: Spring 2023 Problem 2b}

\begin{tcolorbox}[colback=blue!5!white,colframe=blue!75!black,title=Problem]
Previously, he consumed $(4,2)$ when prices were $(1,3)$. Prices are now $(3,5)$. Which bundles $(x_1, x_2)$ would violate WARP if purchased now?
\end{tcolorbox}

\textbf{Solution Strategy:}

\textbf{Step 1:} Calculate the old budget
\[
m^{old} = 1(4) + 3(2) = 10
\]

\textbf{Step 2:} Calculate what the old bundle costs at new prices
\[
\text{Cost}^{new}_{(4,2)} = 3(4) + 5(2) = 22
\]

\textbf{Step 3:} Write the two inequalities
\begin{align*}
\text{Inequality 1:} \quad & x_1 + 3x_2 \leq 10 \quad \text{(was affordable)}\\
\text{Inequality 2:} \quad & 3x_1 + 5x_2 \geq 22 \quad \text{(costs enough now)}
\end{align*}

\textbf{Step 4:} Any bundle satisfying BOTH inequalities violates WARP!

\textbf{Examples:}
\begin{itemize}
\item $(6, 1)$: Old cost = $9 \leq 10$ $eckmark$, New cost = $23 \geq 22$ $eckmark$ $\Rightarrow$ \textbf{VIOLATION}
\item $(3, 2)$: Old cost = $9 \leq 10$ $eckmark$, New cost = $19 \not\geq 22$ $	imes$ $\Rightarrow$ \textbf{SAFE}
\end{itemize}

\newpage

\section{Type 2: Given Cross-Price Information (Multiple Choice)}

\subsection{Format}

You're told:
\begin{itemize}
\item What bundle A cost on Monday and what it cost on Tuesday
\item What bundle B cost on Monday and what it cost on Tuesday
\item What choice was made each day
\end{itemize}

\textbf{Question:} Under what condition is WARP violated?

\subsection{Example: Spring 2024 Q9}

\begin{tcolorbox}[colback=blue!5!white,colframe=blue!75!black,title=Problem]
On Monday, Homer buys a bundle for \$100. On Tuesday, he buys a different bundle for \$80, which would have cost \$100 on Monday. In which situation does Homer violate WARP?

\textbf{(a)} If Monday bundle would cost more than \$80 at Tuesday prices\\
\textbf{(b)} If Monday bundle would cost less than \$80 at Tuesday prices\\
\textbf{(c)} If Monday bundle would cost less than \$100 at Tuesday prices\\
\textbf{(d)} Homer does not violate WARP\\
\textbf{(e)} Not enough information
\end{tcolorbox}

\textbf{Solution Strategy:}

\textbf{Step 1:} Identify what we know
\begin{itemize}
\item Monday: Bought bundle M for \$100
\item Tuesday: Bought bundle T for \$80
\item Bundle T would have cost \$100 on Monday
\end{itemize}

\textbf{Step 2:} Check the two conditions

\textbf{Condition 1:} Was bundle T affordable on Monday?
\begin{itemize}
\item Bundle T cost \$100 on Monday
\item Homer spent \$100 on Monday (his budget)
\item So YES, bundle T was affordable on Monday $eckmark$
\end{itemize}

\textbf{Condition 2:} Was bundle M affordable on Tuesday?
\begin{itemize}
\item Homer bought bundle T for \$80 on Tuesday
\item This proves his budget was at least \$80
\item Bundle M is affordable if it costs $\leq \$80$ at Tuesday prices
\end{itemize}

\textbf{Step 3:} Determine violation condition

For WARP violation, we need BOTH conditions:
\begin{enumerate}
\item Bundle T was affordable Monday $eckmark$ (already true)
\item Bundle M is affordable Tuesday $\Leftrightarrow$ costs $\leq \$80$ at Tuesday prices
\end{enumerate}

\textbf{Answer: (b)} If Monday bundle would cost less than \$80 at Tuesday prices

\subsection{Why This Works}

\textbf{Monday:} He chose M over T when both cost \$100 (both affordable) $\Rightarrow$ M $\succ$ T

\textbf{Tuesday:} If M costs $< \$80$, then both M and T are affordable (he has $\geq \$80$), but he chose T $\Rightarrow$ T $\succ$ M

\textbf{Contradiction!}

\newpage

\section{Type 3: Spring 2025 Version (Slight Variation)}

\begin{tcolorbox}[colback=blue!5!white,colframe=blue!75!black,title=Problem]
On Monday, Homer buys a bundle for \$90. On Tuesday, he buys a different bundle for \$100, that would have cost \$80 at Monday prices. In which situation does he violate WARP?

\textbf{(a)} If Monday bundle would cost \$100 at Tuesday prices\\
\textbf{(b)} If Monday bundle would cost more than \$100 at Tuesday prices\\
\textbf{(c)} Homer does not violate WARP
\end{tcolorbox}

\textbf{Solution:}

\textbf{Given:}
\begin{itemize}
\item Monday: Bought bundle M for \$90 (so budget = \$90)
\item Tuesday: Bought bundle T for \$100 (so budget = \$100)
\item Bundle T would have cost \$80 on Monday
\end{itemize}

\textbf{Check Condition 1:} Was T affordable Monday?
\begin{itemize}
\item T cost \$80 on Monday
\item Budget was \$90
\item \$80 $\leq$ \$90 $eckmark$ YES, T was affordable!
\end{itemize}

\textbf{Check Condition 2:} Was M affordable Tuesday?
\begin{itemize}
\item Homer spent \$100 Tuesday (his budget)
\item M is affordable if it costs $\leq \$100$ at Tuesday prices
\end{itemize}

\textbf{For WARP violation:}
\begin{itemize}
\item Condition 1: $eckmark$ (T was affordable Monday)
\item Condition 2: M costs $\leq \$100$ at Tuesday prices
\end{itemize}

\textbf{Answer: (a)} If Monday bundle would cost \$100 at Tuesday prices (or less!)

Actually both (a) and anything $\leq \$100$ would work, but (b) says "more than \$100" which would NOT violate WARP because then M wouldn't be affordable Tuesday.

\newpage

\section{Step-by-Step Solution Template}

Use this template for ANY WARP problem:

\subsection{Template}

\begin{tcolorbox}[colback=yellow!10!white,colframe=orange!75!black,title=WARP Solution Template]

\textbf{Given Information:}
\begin{itemize}
\item Time 1: Bought bundle $(x_1^*, x_2^*)$ with prices $(p_1^{old}, p_2^{old})$
\item Budget/Cost at Time 1: $m^{old}$
\item Time 2: Considering bundle $(x_1, x_2)$ with prices $(p_1^{new}, p_2^{new})$
\end{itemize}

\textbf{Step 1: Calculate old bundle cost at new prices}
\[
C^{new}_{old} = p_1^{new} x_1^* + p_2^{new} x_2^*
\]

\textbf{Step 2: Write Condition 1 (Was new bundle affordable before?)}
\[
p_1^{old} x_1 + p_2^{old} x_2 \leq m^{old}
\]

\textbf{Step 3: Write Condition 2 (Is old bundle affordable now?)}
\[
p_1^{new} x_1 + p_2^{new} x_2 \geq C^{new}_{old}
\]

\textbf{Step 4: Check both conditions}
\begin{itemize}
\item If BOTH true $\Rightarrow$ \textbf{WARP VIOLATION}
\item If either false $\Rightarrow$ \textbf{NO VIOLATION (safe)}
\end{itemize}

\end{tcolorbox}

\newpage

\section{Common Tricks and Pitfalls}

\subsection{Trick 1: "Would have cost" statements}

When a problem says "bundle X would have cost \$Y at time 2 prices," this is giving you DIRECT information about the cost. Use it!

\textbf{Example:} "Tuesday bundle would have cost \$100 at Monday prices"
\[
\Rightarrow p_1^{Mon} x_1^{Tue} + p_2^{Mon} x_2^{Tue} = 100
\]

\subsection{Trick 2: Budget vs. Cost}

\begin{itemize}
\item If someone \textbf{buys} a bundle for \$X, their budget is \textbf{at least} \$X
\item If we know they spent exactly their whole budget (common assumption), budget = cost
\end{itemize}

\subsection{Trick 3: Inequality direction}

\textbf{ALWAYS:}
\begin{itemize}
\item Past affordability: $\leq$ (bundle was cheap enough)
\item Present affordability: $\geq$ (bundle costs enough to prove old one affordable)
\end{itemize}

\textbf{Never flip these!}

\subsection{Trick 4: "Exactly at the boundary"}

If a bundle costs \textbf{exactly} the budget amount (equality), it's still affordable!
\begin{itemize}
\item \$100 $\leq$ \$100 $eckmark$ TRUE
\item This is a boundary case
\end{itemize}

\newpage

\section{Practice Problems}

\subsection{Problem 1}

Last week: Bought $(3, 5)$ at prices $(2, 4)$\\
This week: Prices are $(5, 3)$

Which bundles violate WARP if purchased this week?
\begin{itemize}
\item (a) $(6, 2)$
\item (b) $(2, 5)$
\item (c) $(4, 4)$
\end{itemize}

\subsection{Problem 2}

Monday: Homer buys bundle M for \$50\\
Tuesday: Homer buys bundle T for \$60, which would have cost \$40 on Monday

WARP is violated if:
\begin{itemize}
\item (a) Bundle M costs less than \$60 at Tuesday prices
\item (b) Bundle M costs more than \$60 at Tuesday prices
\item (c) Never violated
\end{itemize}

\subsection{Problem 3}

Day 1: Prices $(1, 2)$, bought $(10, 5)$\\
Day 2: Prices $(3, 1)$, bought $(4, 8)$

Does this violate WARP? Show your work.

\newpage

\section{Solutions to Practice Problems}

\subsection{Solution 1}

\textbf{Given:}
\begin{itemize}
\item Old: $(3,5)$ at prices $(2, 4)$, cost = $2(3) + 4(5) = 26$
\item New prices: $(5, 3)$
\item Old bundle now costs: $5(3) + 3(5) = 30$
\end{itemize}

\textbf{Conditions:}
\begin{align*}
\text{Condition 1:} & \quad 2x_1 + 4x_2 \leq 26\\
\text{Condition 2:} & \quad 5x_1 + 3x_2 \geq 30
\end{align*}

\textbf{Check each:}

\textbf{(a) $(6, 2)$:}
\begin{itemize}
\item Old cost: $2(6) + 4(2) = 20 \leq 26$ $eckmark$
\item New cost: $5(6) + 3(2) = 36 \geq 30$ $eckmark$
\item \textbf{VIOLATES WARP!}
\end{itemize}

\textbf{(b) $(2, 5)$:}
\begin{itemize}
\item Old cost: $2(2) + 4(5) = 24 \leq 26$ $eckmark$
\item New cost: $5(2) + 3(5) = 25 \not\geq 30$ $	imes$
\item \textbf{SAFE}
\end{itemize}

\textbf{(c) $(4, 4)$:}
\begin{itemize}
\item Old cost: $2(4) + 4(4) = 24 \leq 26$ $eckmark$
\item New cost: $5(4) + 3(4) = 32 \geq 30$ $eckmark$
\item \textbf{VIOLATES WARP!}
\end{itemize}

\subsection{Solution 2}

\textbf{Given:}
\begin{itemize}
\item Monday: M costs \$50 (budget = \$50)
\item Tuesday: T costs \$60 (budget = \$60)
\item T would have cost \$40 on Monday
\end{itemize}

\textbf{Condition 1:} Was T affordable Monday?
\begin{itemize}
\item T cost \$40, budget was \$50
\item \$40 $\leq$ \$50 $eckmark$
\end{itemize}

\textbf{Condition 2:} Is M affordable Tuesday?
\begin{itemize}
\item Budget Tuesday is \$60
\item M is affordable if it costs $\leq \$60$
\end{itemize}

\textbf{Answer: (a)} If M costs less than \$60 at Tuesday prices (or equal!)

\subsection{Solution 3}

\textbf{Day 1:}
\begin{itemize}
\item Bought: $(10, 5)$
\item Prices: $(1, 2)$
\item Cost: $1(10) + 2(5) = 20$
\end{itemize}

\textbf{Day 2:}
\begin{itemize}
\item Bought: $(4, 8)$
\item Prices: $(3, 1)$
\item Cost: $3(4) + 1(8) = 20$
\end{itemize}

\textbf{Check if Day 2 bundle was affordable Day 1:}
\[
1(4) + 2(8) = 20 \leq 20 \quad \checkmark
\]

\textbf{Check if Day 1 bundle is affordable Day 2:}
\[
3(10) + 1(5) = 35
\]

Day 2 budget is \$20 (what they spent). Is \$35 $\leq$ \$20? NO!

\textbf{Conclusion:} \textbf{NO VIOLATION} - Day 1 bundle is NOT affordable on Day 2.

\end{document}
